\documentclass[11pt]{article}

\usepackage[T1]{fontenc}
\usepackage[utf8]{inputenc}
\usepackage{lmodern}
\usepackage[margin=0.92in]{geometry}
\usepackage{microtype}
\usepackage{amsmath,amssymb,amsthm,mathtools}
\usepackage{booktabs,longtable,array,tabularx}
\usepackage{enumitem}
\usepackage{xcolor}
\usepackage{listings}
\usepackage{fancyhdr}
\usepackage{hyperref}
\usepackage{cleveref}

\hypersetup{
  colorlinks=true,
  linkcolor=blue!45!black,
  citecolor=blue!45!black,
  urlcolor=blue!45!black,
  pdftitle={A Kernel-Facing Finite-Use Framework for Mathematics in iPRA},
  pdfauthor={OpenAI, prepared for Vladimir Reshetnikov}
}

\pagestyle{fancy}
\fancyhf{}
\lhead{Kernel-facing finite-use mathematics in iPRA}
\rhead{15 November 2026}
\cfoot{\thepage}
\setlength{\headheight}{14pt}
\setlength{\parindent}{0pt}
\setlength{\parskip}{0.45em}
\setlist{nosep,leftmargin=2em}

\definecolor{codebg}{RGB}{246,247,249}
\definecolor{codekw}{RGB}{30,70,130}
\lstset{
  basicstyle=\ttfamily\small,
  backgroundcolor=\color{codebg},
  frame=single,
  rulecolor=\color{black!15},
  breaklines=true,
  columns=fullflexible,
  keepspaces=true,
  keywordstyle=\color{codekw}\bfseries,
  showstringspaces=false
}

\newtheorem{theorem}{Theorem}[section]
\newtheorem{lemma}[theorem]{Lemma}
\newtheorem{proposition}[theorem]{Proposition}
\newtheorem{corollary}[theorem]{Corollary}
\newtheorem{definition}[theorem]{Definition}
\newtheorem{remark}[theorem]{Remark}
\newtheorem{example}[theorem]{Example}
\newtheorem{warning}[theorem]{Warning}

\newcommand{\N}{\mathbb N}
\newcommand{\Z}{\mathbb Z}
\newcommand{\Q}{\mathbb Q}
\newcommand{\B}{\{0,1\}}
\newcommand{\PRA}{\mathsf{iPRA}}
\newcommand{\PRAAx}{\mathsf{PRAAx}}
\newcommand{\Prf}{\mathsf{Prf}}
\newcommand{\QF}{\mathsf{QF}}
\newcommand{\Code}{\mathsf{Code}}
\newcommand{\True}{\mathsf{true}}
\newcommand{\False}{\mathsf{false}}
\newcommand{\dist}{\mathop{\mathrm{dist}}}
\newcommand{\obs}{\mathop{\mathrm{obs}}}
\newcommand{\valid}{\mathop{\mathrm{Valid}}}
\newcommand{\viable}{\mathop{\mathrm{Viable}}}
\newcommand{\pair}[2]{\langle #1,#2\rangle}
\newcommand{\sequent}[3]{#1\vdash_{#2}#3}

\title{\bfseries A Kernel-Facing Finite-Use Framework\\
for Mathematics in Intuitionistic Primitive Recursive Arithmetic\\[0.4em]
\large Proof-carrying DAGs, bounded relational choice, stable observables,\\
and explicit boundaries for difficult mathematics}
\author{Prepared for Vladimir Reshetnikov}
\date{15 November 2026}

\begin{document}
\maketitle

\begin{abstract}
The current \texttt{Logic/PrimitiveRecursiveArithmetic} tree of the ProveIt
repository contains a kernel-checked foundation and a large collection of
research packets spanning algebra, analysis, geometry, topology, number
theory, probability, category theory, and mathematical physics.  The packets
share a productive finitary methodology---codes, bounded traces, finite
certificates, rational approximations, and explicit resource budgets---but
many stop one level before the formal target: they establish a host-language
fact or design a semantic algorithm without constructing a derivation of the
repository's deep-embedded judgment
\[
  \Gamma\vdash_{\PRAAx}\varphi.
\]
A recent contract-DAG packet identifies the right architecture but leaves its
generic object compiler schematic.

This report develops a kernel-facing finite-use framework that closes that
generic logical layer.  A proof-carrying reverse-topological schedule records,
for each node of a fixed finite contract graph, an actual local
\(\Prf\,\PRAAx\) derivation from earlier nodes.  Simultaneous cut compiles a new
node, and finite iteration compiles the whole graph.  A separate
quantifier-free loop theorem packages the sole induction needed for bounded
traces while permitting a quantified terminal result.  Bounded search and
backward viability eliminate hidden finite choice; stable-observable
comparison eliminates the need for coherent representatives; exact branch
selection is permitted only under an explicit gap; and genuinely global or
ineffective mathematics is isolated as a boundary theorem rather than
misrepresented as an algorithm.

The main arguments are given in full at the human level.  The archive includes
Lean modules whose generic conclusions have the actual embedded
\texttt{Prf PRAAx} type, a dependency-ordered formalization plan, a conservative
Lean schedule generator, exact Python reference algorithms, and independently
checked certificates.  Lean was unavailable in the build runtime, so source
without placeholders is supplied but not claimed to have been typechecked.
The framework does not manufacture missing height bounds, convergence moduli,
separation estimates, or effective forms of deep global theorems; it makes
those obligations precise and composable.
\end{abstract}

\tableofcontents
\newpage

\section{Executive summary}

The useful question is not whether all mathematics can be encoded by natural
numbers.  That is routine.  The useful question is:

\begin{quote}
What finite data, resource transformer, verifier, and object-language proof are
sufficient for the particular observation demanded by a later theorem?
\end{quote}

The answer developed here is a six-layer compiler.

\begin{enumerate}
\item \textbf{Fixed coding.}  Select concrete primitive-recursive codes for the
      objects used by one theorem family.  Do not postulate a universal total
      interpreter for arbitrary codes.
\item \textbf{Backward demand.}  Starting from the final finite observation,
      compute the precision, fuel, height, support, and branch demands required
      by predecessor nodes.
\item \textbf{Finite production.}  Run a bounded trace, bounded search, or
      finite relational computation.  Producers may return a cloud of valid
      witnesses instead of a single coherent representative.
\item \textbf{Local verification.}  Each node supplies a quantifier-free graph
      predicate and actual embedded base/step/finish derivations.
\item \textbf{Kernel composition.}  A reverse-topological proof schedule
      composes local derivations by simultaneous cut, producing
      \(\Prf\,\PRAAx\) theorems for all nodes.
\item \textbf{Erasure.}  Stable comparison, a modulus, an isolation gap, or an
      explicit boundary theorem converts certificate data into the requested
      mathematical statement.
\end{enumerate}

The central advance is layer 5.  Earlier packets already contain excellent
instances of layers 1--4 and 6, and a few isolated object-level combinators.
The new \texttt{ProofBuild} construction turns those pieces into a genuine
finite DAG compiler without weakening the formal target to \texttt{True}.

\paragraph{Resolved at the generic level.}
Finite DAG composition, bounded loop reflection, bounded relational choice,
stable witness erasure, and gap-guarded branch selection all receive complete
human proofs and kernel-facing Lean statements.

\paragraph{Not resolved by abstraction alone.}
A framework cannot create a missing effective Siegel bound, Faltings height
bound, spectral gap, root-separation estimate, convergence modulus, or
prime-distribution constant.  Such facts appear as explicit producer or
boundary fields.  This is a feature: the formal statement now says exactly
where the deep mathematics is used.

\section{Repository baseline and diagnosis}
\label{sec:audit}

\subsection{The formal target}

The audited foundation is \texttt{ConciseIPRA}.  Its ordinary presentation of
intuitionistic primitive recursive arithmetic has:

\begin{itemize}
\item a type \texttt{PR k} of primitive-recursive function symbols of arity
      \(k\);
\item scoped terms \texttt{PTm n};
\item scoped formulas \texttt{Fm PTm n};
\item an inductive derivability relation \texttt{Prf PRAAx Gamma phi};
\item defining equations for all fixed primitive-recursive symbols,
      successor axioms, and induction for quantifier-free formulas;
\item substitution, weakening, simultaneous cut, equality congruence, and a
      bridge to the native scalar calculus.
\end{itemize}

Thus the formal target is not a host proposition of the form
\texttt{Algorithm x = y}.  It is a term inhabiting
\[
  \Prf(\PRAAx,\Gamma,\varphi),
\]
written in the repository as
\texttt{Gamma \(\vdash\)[PRAAx] phi}.  A semantic Lean proof may be useful
for designing an algorithm, but it is not automatically a theorem of iPRA.
The conversion must be exhibited through the deep-embedded deduction rules.

\subsection{The current research corpus}

At the audited revision, the proof index describes 72 independent theorem or
framework packets and explicitly warns that they are not thereby verified by
the maintained foundation.  The foundation subtable contains 17 general
frameworks.  Several of the newest are highly relevant:

\begin{itemize}
\item finite observations, finite-use relations, evidence-indexed witnesses,
      and stable observables;
\item changing-parameter traces, resource vectors, lexicographic descent, and
      bounded enumeration;
\item coherent finite presentations, contract DAGs, and categorical finite
      presentations;
\item object-level loop and observation combinators, and small categorical
      calculations concluding \texttt{Prf PRAAx}.
\end{itemize}

The recurring final gap is accurately summarized by the recent
\texttt{ipra\_cfp\_contract\_compiler\_2026-11-15} packet: the architecture of
reverse demand, forward trace, and local proof-producing nodes is present, but
the generic object-language DAG theorem remains schematic.  In its Lean
blueprint, the intended theorem still concludes \texttt{True}.  In contrast,
\texttt{ObjectObservation.lean} and \texttt{ObjectCategory.lean} demonstrate
that small, honest \texttt{Prf PRAAx} combinators can be written against the
current API.

\subsection{Twelve recurring difficulties}

\begin{longtable}{p{0.06\textwidth}p{0.35\textwidth}p{0.50\textwidth}}
\toprule
ID & Difficulty & Required correction \\
\midrule
\endhead
D1 & Host truth versus object proof & Require every formal claim to state
whether it is a host theorem, a verified certificate, or an embedded
\(\Prf\,\PRAAx\) derivation. \\
D2 & Missing generic DAG compiler & Compile a fixed topological schedule by
simultaneous cut, not by an object-level universal interpreter. \\
D3 & Quantifier-free induction restriction & Put all changing data into one
coded state and use one QF invariant; allow quantified extraction only after
the finite loop. \\
D4 & Hidden coherent choice & Compute finite viability tables and bounded least
witnesses, or retain a relation when no canonical witness is needed. \\
D5 & Approximation collapsed to exactness & Demand a separate gap, apartness,
uniqueness radius, or rounding margin. \\
D6 & Ineffectivity disguised as computation & Make the global theorem an
explicit boundary premise. \\
D7 & Universal evaluator assumption & Specialize every fixed program before
formalization; a total PR universal evaluator cannot exist. \\
D8 & Changing parameters in recursion & Store the current parameter and
resource counters in the state, with fixed PR transition. \\
D9 & Cardinality confused with enumeration & Require a code-height cutoff and
a decidable validity predicate, not merely a finite-cardinality theorem. \\
D10 & Quotients and tie-breaking & Compare observables from any two valid
representatives; erase the representative only after stability is proved. \\
D11 & Convergence without a rate & Use explicit moduli, backward precision
demands, or finite metastable windows. \\
D12 & Deep domain estimates absent & Separate reusable finite algebra from
missing global estimates and record the latter as producer/boundary fields. \\
\bottomrule
\end{longtable}

The framework below resolves D1--D5 and D7--D8 at the reusable level, clarifies
D9--D11, and gives a disciplined interface for D6 and D12.

\section{Design principles}

\subsection{Fixed families, not universal interpretation}

Every theorem packet uses finitely many concrete algorithms.  Their PR syntax
can be generated outside the object theory and then inserted as fixed function
symbols.  The defining equations already belong to \(\PRAAx\).  This is both
simpler and mathematically necessary: a total primitive-recursive evaluator
for all unary primitive-recursive programs would contradict diagonalization;
see \cref{prop:no-universal}.

\subsection{Finite use before global objects}

A downstream proof rarely consumes an infinite object in full.  It asks for a
finite observation: the first \(m\) coefficients, an interval of width
\(2^{-p}\), a finite chain group, a bounded list of candidates, a sign, or a
finite matrix identity.  The framework therefore specifies a demand
\(d\) and proves only what the client uses.

\subsection{Relations before canonical witnesses}

Canonical representatives are expensive and often mathematically irrelevant.
Roots collide; eigenvectors rotate inside an eigenspace; subdivisions admit
many carrier maps; local analytic branches permute.  A finite relation of valid
outputs is usually easier to compute constructively.  A comparison theorem can
then show that all valid outputs have the same requested observation.

\subsection{Exactness is a separate phase}

Approximate data supports approximate conclusions.  Exact equality, exact zero,
or a discrete branch requires a gap.  The gap may be as simple as the unit
spacing of integers or as deep as a root-separation theorem.  It must be stated
and proved separately.

\subsection{Boundary honesty}

Some classical theorems are finitary in conclusion but ineffective in their
known proofs.  The correct iPRA result may therefore be conditional:
\[
  \mathsf{Boundary}(x)\longrightarrow \mathsf{FiniteConclusion}(x).
\]
This is a useful formalization.  It identifies the exact missing ingredient
without pretending to have computed it.

\section{Concrete finite coding and resource discipline}
\label{sec:coding}

The framework is independent of a particular Gödel coding, but an integration
should choose one fixed library and prove its equations once.  An explicit
option is the positive binary pair
\[
  \pair{a}{b}=2^a(2b+1).
\]
For positive \(z\), the left projection is the largest \(a\le z\) such that
\(2^a\mid z\), found by bounded search; the right projection is
\(\bigl(z/2^a-1\bigr)/2\).  Give zero a fixed default decoding.  Exponentiation,
bounded divisibility tests, bounded search, and bounded division are primitive
recursive, hence so are both projections.  Finite tuples and lists are nested
pairs, preferably with a separately stored length.

Signed integers may be encoded by parity,
\[
  \mathsf{int}(z)=
  \begin{cases}
  2z,&z\ge 0,\\
  -2z-1,&z<0,
  \end{cases}
\]
and a rational by a signed numerator and a positive denominator.  Normalization
is useful but not always necessary: equality and order can be checked by
cross-multiplication.  Exact rational arithmetic is therefore PR and its graph
is quantifier-free when expressed as equality between fixed PR terms.

\begin{definition}[Resource vector]
A resource vector is a finite tuple
\[
  \rho=(p,f,h,s,b,\ldots)
\]
of requested precision, loop fuel, code-height cutoff, support bound, branch
budget, and any domain-specific counters.  A node demand transformer is a fixed
PR map from the requested output resource to the resources required of its
parents.
\end{definition}

For a fixed acyclic graph, reverse traversal and coordinatewise maximum compute
all demands.  This traversal is performed at code-generation time when the
graph itself is fixed.  If the graph shape is input data, it is instead a
bounded trace with a QF invariant, not a universal meta-interpreter.

\begin{lemma}[Closure of fixed bounded data operations]
For a fixed coding library, bounded list lookup, update, append, Boolean folds,
bounded minimum, finite-table reversal, and reverse-demand propagation are
primitive recursive.
\end{lemma}

\begin{proof}
Each operation is a composition of projections, fixed arithmetic, and a finite
primitive recursion whose recursion argument is the supplied bound or list
length.  Lookup scans at most the stored length; update rebuilds the same
bounded prefix; a Boolean fold carries one accumulator bit; bounded minimum
carries a pair consisting of a found flag and the least witness seen so far.
Reverse propagation carries a finite coded table and performs one bounded
update per edge.  Closure of primitive-recursive functions under composition
and primitive recursion gives the claim.  No unbounded minimization occurs.
\end{proof}

\section{Finite-use contracts}

\begin{definition}[Operational contract]
A finite-use operational contract for a node \(v\) consists of:
\begin{enumerate}
\item a premise formula \(P_v\) and result formula \(A_v\);
\item a finite list of predecessor nodes \(\operatorname{par}(v)\);
\item fixed PR codes for demand, production, and verification;
\item a QF loop invariant when production is iterative;
\item an embedded local derivation
\[
  (A_u)_{u\in\operatorname{par}(v)},P_v,\Gamma
  \vdash_{\PRAAx} A_v;
\]
\item optional comparison, gap, and boundary derivations.
\end{enumerate}
\end{definition}

The local derivation is the decisive field.  A Python checker or a theorem
about host naturals is useful evidence but does not fill this field.  Conversely,
the local derivation need not reprove generic logic: the kernel compiler below
handles all cuts and context management.

\begin{definition}[Reverse-topological proof build]
Fix a base context \(\Gamma\).  A proof build for a list
\(D=[A_{k-1},\ldots,A_0]\) is the assertion
\[
  \forall A\in D,\quad \Gamma\vdash_{\PRAAx}A.
\]
The list is reverse topological: newly compiled nodes are consed to its front.
\end{definition}

This definition is deliberately a metalevel Lean structure, not an arithmetic
encoding inside iPRA.  Each theorem packet has a fixed finite dependency graph;
Lean generates and checks the corresponding proof term.  Variable-size traces
inside a theorem are handled by QF induction, not by asking iPRA to interpret
arbitrary Lean graphs.

\section{Kernel-facing proof-DAG compilation}
\label{sec:dag}

\begin{theorem}[Proof-carrying step]
\label{thm:step}
Let \(D\) be an existing proof build over \(\Gamma\), and suppose
\[
  D,\Gamma\vdash_{\PRAAx} A.
\]
Then \([A]D\) is a proof build over \(\Gamma\).
\end{theorem}

\begin{proof}
Every old member of \(D\) retains its stored derivation.  For the new formula
\(A\), apply simultaneous cut to the local derivation.  If a used hypothesis
belongs to \(D\), replace it by the derivation stored in the old build.  If it
belongs to \(\Gamma\), replace it by the hypothesis rule.  The result is
\(\Gamma\vdash_{\PRAAx}A\).  All operations are constructors of the
inductive relation \(\Prf\); no semantic soundness theorem is invoked.
\end{proof}

The supplied Lean definition is essentially the proof above:

\begin{lstlisting}[language={}]
def ProofBuild.step (s : ProofBuild Gamma done)
    (local : done ++ Gamma |- [PRAAx] A) :
    ProofBuild Gamma (A :: done) where
  derive := by
    intro C hC
    rcases List.mem_cons.mp hC with rfl | hOld
    * exact local.trans (fun D hD => by
        rcases List.mem_append.mp hD with hDone | hBase
        * exact s.derive D hDone
        * exact Prf.hyp hBase)
    * exact s.derive C hOld
\end{lstlisting}

Here the displayed ASCII \texttt{|-} abbreviates the repository's notation.
The actual file uses the Unicode notation imported from
\texttt{ConciseIPRA.Deduction}.

\begin{theorem}[Dependency-restricted step]
\label{thm:step-from}
If \(E\subseteq D\) and
\(E,\Gamma\vdash_{\PRAAx}A\), then \([A]D\) is a proof build.
\end{theorem}

\begin{proof}
Apply the same simultaneous cut, using the inclusion certificate to retrieve
stored derivations only for hypotheses in \(E\).  This is
\texttt{ProofBuild.stepFrom}.  The inclusion proof is the metalevel witness
that all dependencies precede the new node.
\end{proof}

\begin{theorem}[Finite DAG compilation]
\label{thm:dag}
Let \(G\) be a finite acyclic graph with a fixed topological order
\(v_0,\ldots,v_{m-1}\).  Associate a formula \(A_i\) to each node and suppose
that for each \(i\) there is an embedded local derivation of \(A_i\) from
\(\Gamma\) and formulas attached to predecessors of \(v_i\).  Then for every
\(i<m\),
\[
  \Gamma\vdash_{\PRAAx}A_i.
\]
\end{theorem}

\begin{proof}
Start with the empty proof build.  At stage \(i\), all predecessor formulas are
already in the reverse-topological list.  Apply
\cref{thm:step-from}.  Finite induction on \(i\) constructs a proof build
containing all nodes.  Retrieval of any member gives the desired embedded
derivation.
\end{proof}

\begin{corollary}[Untrusted generation, trusted checking]
A program may generate the sequence of \texttt{ProofBuild.step} applications
from a JSON dependency graph.  Correctness of the generator is not part of the
trusted base: Lean checks the emitted term against the local derivation types.
\end{corollary}

This is the correct role for code generation.  It saves bookkeeping without
introducing a universal object-level evaluator or a new axiom.

\section{Quantifier-free loop compilation}
\label{sec:loop}

The DAG compiler composes local theorems.  We now show how a bounded iterative
producer supplies one local theorem using the sole induction schema available
in iPRA.

Let \(P,Q\) be formulas over \(n\) free variables, let
\(I\) be a formula over \(n+1\) variables whose index-0 variable is the loop
stage, and let \(t\) be a term for the fuel.  Assume:
\begin{align*}
  & I\text{ is quantifier-free};\\
  & P,\Gamma\vdash_{\PRAAx} I(0);\\
  & I(i),P,\Gamma\vdash_{\PRAAx}I(i+1);\\
  & I(t),P,\Gamma\vdash_{\PRAAx}Q.
\end{align*}
The exact scoped versions shift \(P\) and \(\Gamma\) in the induction step, as
implemented in the Lean source.

\begin{theorem}[Loop reflection under a precondition]
\label{thm:loop}
Under the four assumptions above,
\[
  \Gamma\vdash_{\PRAAx}P\to Q.
\]
\end{theorem}

\begin{proof}
Apply the iPRA QF-induction rule to the base and step derivations.  This yields
\(P,\Gamma\vdash\forall i\,I(i)\).  Universal elimination at the fuel term
\(t\) gives \(P,\Gamma\vdash I(t)\).  Cut this into the finish derivation to
obtain \(P,\Gamma\vdash Q\).  Finally introduce the implication.
\end{proof}

The terminal formula \(Q\) may contain quantifiers.  Only the loop invariant
must be QF.  This separation is crucial: bounded arithmetic facts are proved
inside the invariant, then a concrete witness or finite code is exported by
existential introduction.

\subsection{Changing parameters}

Suppose an algorithm changes a precision parameter, chart index, degree cap,
or support bound during its run.  There is no need for a stronger induction
schema.  Encode the entire configuration
\[
  s_i=\pair{i}{\pair{p_i}{\pair{\rho_i}{\text{domain state}_i}}}
\]
as one natural.  The transition \(s_{i+1}=\mathsf{step}(s_i,x)\) is a fixed PR
function.  The invariant states, using fixed Boolean checkers, that:

\begin{itemize}
\item the stage component equals \(i\);
\item all local transitions up to \(i\) passed their checkers;
\item every resource component is within the reverse-computed budget;
\item any lexicographic descent measure has decreased or the state is frozen.
\end{itemize}

Because bounded universal assertions are implemented by PR Boolean folds, the
invariant is an equality of terms and hence QF.  A primitive-recursive fuel
must be supplied.  General well-founded recursion is not smuggled in: the
resource proof explicitly bounds the number of transitions.

\section{Bounded search and finite relational choice}
\label{sec:choice}

\subsection{Least bounded witness}

Let \(b(x,y)\in\B\) and \(B(x)\in\N\) be fixed PR functions.  Define a scan
state \((f,w)\), where \(f\) is a found flag.  Initially \((0,0)\).  At stage
\(i\le B(x)\), keep the state if \(f=1\); otherwise set \((1,i)\) when
\(b(x,i)=1\), and keep \((0,0)\) when it is zero.

\begin{theorem}[Bounded least search]
\label{thm:bounded-search}
The scan is primitive recursive.  After \(B(x)+1\) stages:
\begin{enumerate}
\item the flag is one iff some \(y\le B(x)\) satisfies \(b(x,y)=1\);
\item when the flag is one, the stored \(w\) is the least such \(y\).
\end{enumerate}
Both assertions admit QF-induction proofs after bounded quantifiers are
compiled to Boolean folds.
\end{theorem}

\begin{proof}
Primitive recursiveness follows from the explicit finite-state transition.
For correctness, use the stage invariant:
\[
\begin{split}
 f_i=0 &\Longleftrightarrow \bigwedge_{y<i}(b(x,y)=0),\\
 f_i=1 &\Longrightarrow b(x,w_i)=1\ \wedge\ w_i<i\ \wedge
       \bigwedge_{y<w_i}(b(x,y)=0).
\end{split}
\]
Replace both bounded conjunctions by fixed PR Boolean folds.  The base case is
immediate.  In the step case, if a witness was already found the state is
frozen; otherwise the value of \(b(x,i)\) gives the two branches and preserves
the displayed clauses.  At stage \(B(x)+1\) the first equivalence proves
existence detection, and the second clause proves minimality.
\end{proof}

This theorem eliminates finite choice only when the bound is explicit.  An
unbounded existential does not become constructive merely because a classical
proof says its solution set is finite.

\subsection{Backward viability}

Let layers \(L_0,\ldots,L_m\) be finite coded candidate tables.  Let
\(C_i(x)\) say that \(x\) is a valid candidate in layer \(i\), and let
\(E_i(x,y)\) be a decidable compatibility relation from layer \(i\) to
\(i+1\).  Define viability backward:
\begin{align*}
 V_m(x)&=C_m(x),\\
 V_i(x)&=C_i(x)\wedge
  \bigvee_{y<N_{i+1}}\bigl(E_i(x,y)\wedge V_{i+1}(y)\bigr).
\end{align*}
The disjunction is a PR Boolean fold over the explicit next-layer bound.

\begin{theorem}[Viability correctness]
\label{thm:viability}
For every \(i\le m\), \(V_i(x)=1\) iff \(x\) begins a compatible path through
layers \(i,i+1,\ldots,m\).
\end{theorem}

\begin{proof}
Proceed by backward induction on \(m-i\).  At \(i=m\), a path consists of the
single candidate and the statement is the definition of \(V_m\).  For the
step, if \(V_i(x)=1\), bounded least search yields a concrete \(y<N_{i+1}\)
with \(E_i(x,y)=1\) and \(V_{i+1}(y)=1\).  The induction hypothesis extends
\(y\) to a path, so prefixing \(x\) gives a path from layer \(i\).  Conversely,
if a path begins with \(x,y\), then \(C_i(x)=1\), \(E_i(x,y)=1\), and the tail
makes \(V_{i+1}(y)=1\) by induction.  Hence the finite disjunction in the
definition of \(V_i\) is true.
\end{proof}

\begin{corollary}[Lexicographic viable path]
\label{cor:lex-path}
If a full path exists, repeatedly choosing the least viable compatible
successor constructs a full path by a bounded primitive recursion.  No form of
countable or dependent choice is used.
\end{corollary}

The proof is constructive because the viability table gives an explicit
nonempty finite search domain at every step.  In applications where no single
path is needed, one may stop earlier and reason directly with the relation.

\section{Stable observables and guarded exactness}
\label{sec:stable}

\subsection{Witness-independent conclusions}

\begin{theorem}[Stable-observable erasure]
\label{thm:stable}
Suppose
\[
  \Gamma\vdash C_1,\qquad \Gamma\vdash C_2,
  \qquad C_1,C_2,\Gamma\vdash O.
\]
Then \(\Gamma\vdash O\).
\end{theorem}

\begin{proof}
Introduce implications twice in the comparison derivation, obtaining
\(\Gamma\vdash C_2\to(C_1\to O)\).  Apply modus ponens first to the derivation
of \(C_2\), then to the derivation of \(C_1\).  This is exactly the supplied
\texttt{stableObservable} proof term.
\end{proof}

Although logically elementary, this theorem changes the architecture of many
formalizations.  Producers at different precisions may return unrelated roots,
eigenvectors, charts, subdivisions, or minimizers.  The comparison theorem,
not a coherent global selection, carries the mathematical burden.

For a rational-valued observation, a common comparison formula is
\[
  \valid(c_1,p)\wedge\valid(c_2,p)
  \longrightarrow
  |\obs(c_1)-\obs(c_2)|\le 2^{-p}.
\]
A canonical name can be obtained by choosing the least valid certificate at
each precision, but its correctness depends only on pairwise stability.  The
chosen certificates need not form a coherent chain.

\subsection{Exactness from a gap}

\begin{theorem}[Gap decision]
\label{thm:gap}
Let \(z\) be known to satisfy the disjunction
\[
  z=0\quad\vee\quad |z|\ge\delta
\]
for a positive rational \(\delta\).  If a rational \(q\) satisfies
\(|q-z|<\delta/3\), then the decidable comparison
\(|q|<\delta/2\) distinguishes the two branches correctly.
\end{theorem}

\begin{proof}
If \(z=0\), then \(|q|<\delta/3<\delta/2\).  If
\(|z|\ge\delta\), the reverse triangle inequality gives
\[
  |q|\ge |z|-|q-z|>\delta-\delta/3=2\delta/3>\delta/2.
\]
Thus the rational threshold test returns different answers in the two cases.
All inequalities are finite rational calculations.
\end{proof}

\begin{warning}
An approximation alone does not decide exact zero.  The disjunction or an
equivalent isolation theorem is an essential hypothesis.  When it is absent,
the correct output remains an interval, a relation, or an apartness-conditional
statement.
\end{warning}

Integer recovery is the special case \(\delta=1\): an approximation within
\(1/2\) uniquely determines an integer.  Root labels require a root-separation
bound; spectral projectors require a spectral gap; a unique minimizer requires
a quantitative uniqueness radius.

\subsection{Setoids and quotients}

Arithmetic equality of codes is often too strong.  A finite presentation may
have a setoid equality \(R(a,b)\), and a topological construction may admit
many equivalent representatives.  The object development must provide
embedded reflexivity, symmetry, transitivity, and operation congruence for
\(R\).  Stable observations should be formulated in terms of \(R\), not Lean's
host quotient.  This keeps all representative erasure visible in the object
proof.

\section{Boundary factorization and genuine obstructions}
\label{sec:boundary}

\begin{theorem}[Boundary factorization]
\label{thm:boundary}
If \(\Gamma\vdash B\) and \(B,\Gamma\vdash G\), then
\(\Gamma\vdash G\).  If the first derivation is unavailable, the honest result
is \(\Gamma\vdash B\to G\).
\end{theorem}

\begin{proof}
Introduce the implication in the second derivation and apply it to the first.
The conditional version is simply implication introduction without the final
application.
\end{proof}

The theorem is trivial logically but decisive methodologically.  It gives a
standard location for a global height theorem, a zero-free region, a compactness
principle with a modulus, or any other ingredient not generated by the finite
producer.

\subsection{Cardinality is not a cutoff}

A statement ``there are at most \(N\) solutions'' does not provide a number
\(H\) such that every solution code is at most \(H\).  Enumeration requires:

\begin{enumerate}
\item a decidable validity checker \(V(x)\);
\item a proved code-height cutoff \(H\);
\item bounded search through \(x\le H\);
\item a completeness proof that every solution appears below \(H\).
\end{enumerate}

Without item 2, finite cardinality cannot be converted into a primitive-
recursive complete list.  This distinction is central for integral points,
rational points, finitely generated groups, and classification statements.

\subsection{No total PR universal evaluator}

\begin{proposition}[Diagonal obstruction]
\label{prop:no-universal}
There is no primitive-recursive function \(U(e,x)\) that evaluates every unary
primitive-recursive function from a code \(e\).
\end{proposition}

\begin{proof}
Assume such \(U\) exists.  Then
\(g(e)=U(e,e)+1\) is primitive recursive, so it has some code \(d\).  Universality
would give \(g(d)=U(d,d)\), while the definition gives
\(g(d)=U(d,d)+1\), a contradiction.
\end{proof}

This obstruction does not prevent syntax from being encoded or bounded
interpreters from being PR.  It prevents the framework from using a total
universal evaluator as a shortcut.  Fixed-program specialization is the right
solution.

\section{Domain adapters}
\label{sec:domains}

This section describes what the six compiler layers mean in several domains.
It is a formalization map, not a claim that every deep theorem is already
proved in iPRA.

\subsection{Algebra and combinatorics}

Finite groups, rings, matrices, graphs, partitions, and tableaux fit directly.
Use bounded table codes; compute normalization or elimination traces; verify
local equations; compile the trace; and export a finite witness.  Smith normal
form, Gaussian elimination over rationals, determinant identities, finite
orbit calculations, and combinatorial recurrences are natural early targets.
The main care is to include explicit dimension and coefficient-height bounds.

\subsection{Analysis and differential equations}

Represent a real or complex quantity by a PR procedure returning rational
approximations with an explicit error contract.  A client request \(p\) is
pushed backward through operations.  Addition uses essentially the same
precision; multiplication also needs magnitude bounds; inversion requires an
apartness lower bound; composition uses the outer modulus to request inner
precision.  For a polynomial ODE, a bounded Picard or Taylor trace can produce
a finite interval certificate.  A global solution theorem requires explicit
existence intervals and growth bounds, recorded as boundary data when not
constructively available.

When a full convergence rate is unavailable but every finite client uses only
a finite window, a metastable contract is often sufficient:
\[
  \forall p,g\ \exists N\le F(p,g)\ \forall i,j\in[N,N+g(N)],
  |a_i-a_j|\le 2^{-p}.
\]
To remain in iPRA, \(F\) must be explicit and the bounded window checker must be
PR.  Metastability does not magically supply an ineffective \(F\).

\subsection{Topology}

For finite simplicial or cubical complexes, encode vertices, simplices,
incidence, chains, and boundary matrices by bounded tables.  Subdivision and
carrier maps are finite traces.  Different tie-breaking rules need not produce
identical maps; it is enough to construct a bounded prism or chain homotopy and
show that the induced homology observation is stable.  Global theorems such as
Jordan separation or dimension theory then split into a finite combinatorial
compiler plus a theorem-specific approximation or separation boundary.

\subsection{Category theory}

A finite category presentation consists of bounded object and arrow tables,
source/target maps, composition, and setoid equality, together with checked
laws.  Limits, representable functors, finite Kan extensions, and finite sheaf
gluing become finite relational contracts.  The existing object-level inverse
uniqueness and Yoneda component calculations can serve as local nodes.  The new
DAG compiler composes such nodes without reimplementing category logic.

\subsection{Algebraic geometry}

Finite algebraic geometry separates cleanly from global arithmetic geometry.
Polynomial evaluation, derivatives, resultants, Macaulay matrices, bounded
row reduction, finite chart covers, and certificates of ideal membership are
finite algebra and belong in trace nodes.  Complete rational or integral point
lists require additional height/cutoff theorems.  Those theorems are producer
fields when effective and boundary fields otherwise.  This prevents a finite
Macaulay calculation from being mistaken for an effective form of Faltings or
Siegel.

\subsection{Number theory}

Congruence checking, finite local solubility, bounded descent, explicit
factorization certificates, and fixed-exponent searches are straightforward.
Analytic distribution theorems require explicit error terms and constants.
Exceptional-zero issues, ineffective constants, or merely existential height
bounds must remain visible.  Conditional finite consequences are still
valuable and can be compiled using \cref{thm:boundary}.

\subsection{Probability and asymptotics}

Finite distributions, exact moments, couplings on bounded sample spaces, and
finite concentration checks are PR.  Limit theorems need a quantitative rate or
a metastable bound.  A central-limit packet should therefore separate the
finite characteristic-function or moment computation from the analytic
uniformity theorem that turns it into a distributional approximation.

\section{Worked example: exact cubic bisection}
\label{sec:cubic}

Consider
\[
  f_x(y)=y^3+y-x
\]
for rational \(x\).  The goal is not to invoke classical real analysis.  For
each requested precision \(p\), we construct a finite rational interval
\([a_p,b_p]\) of width at most \(2^{-p}\) with
\(f_x(a_p)\le0\le f_x(b_p)\).  These intervals are a finite-use name for the
unique root.

\begin{lemma}[Quantitative strict monotonicity]
\label{lem:cubic-monotone}
If \(u<v\) are rational, then
\[
  f_x(v)-f_x(u)\ge v-u>0.
\]
\end{lemma}

\begin{proof}
Factor:
\[
 f_x(v)-f_x(u)=(v-u)(v^2+uv+u^2+1).
\]
The second factor is at least one because
\[
  4(v^2+uv+u^2)=(2u+v)^2+3v^2\ge0.
\]
Multiplying by the positive \(v-u\) gives the result.
\end{proof}

The stronger lower bound implies inverse stability:
\[
  |u-v|\le |f_x(u)-f_x(v)|.
\]
It also proves uniqueness without excluded middle: two exact zeros cannot be
strictly ordered in either direction.

Let
\[
  M=\lceil |x|\rceil+1.
\]
Then \(f_x(-M)<0<f_x(M)\).  Indeed \(M>|x|\), and the cubic terms only strengthen
the inequalities.

At each stage, put \(m=(a+b)/2\).  Rational sign is decidable.  If
\(f_x(m)\le0\), replace \(a\) by \(m\); otherwise replace \(b\) by \(m\).

\begin{lemma}[Bisection invariant]
After every stage,
\[
  a\le b,\qquad f_x(a)\le0\le f_x(b),
\]
and the width is half its previous value.
\end{lemma}

\begin{proof}
The midpoint lies between the endpoints.  In the first branch, the new lower
endpoint has nonpositive value and the old upper endpoint retains nonnegative
value.  In the second branch, the old lower endpoint is retained and the new
upper endpoint has positive value.  In both branches the new interval is one
of the two halves.
\end{proof}

Define
\[
  k=p+\left\lceil\log_2(2M)\right\rceil.
\]
The logarithm here is the least natural \(r\) with \(2M\le2^r\), computed by
bounded search or bit length.

\begin{theorem}[Cubic finite-use certificate]
\label{thm:cubic}
After \(k\) bisection stages, the sign bracket is preserved and
\[
  b-a=\frac{2M}{2^k}\le2^{-p}.
\]
The producer and verifier are primitive recursive on rational codes.
\end{theorem}

\begin{proof}
The invariant follows by finite induction from the preceding lemma.  Repeated
halving gives the exact width formula.  By the definition of the ceiling
logarithm, \(2M\le2^{k-p}\), hence \(2M/2^k\le2^{-p}\).  Every stage uses exact
integer/rational arithmetic and one decidable sign test, so the full trace and
its checker are primitive recursive.
\end{proof}

\subsection{Object-iPRA compilation}

A complete object development needs fixed PR codes for normalized rationals,
addition, multiplication, comparison, ceiling, and ceiling logarithm.  Define
the trace state to contain \((i,a_i,b_i)\).  The QF invariant is a conjunction
of checker equalities asserting:

\begin{enumerate}
\item the state is the result of \(i\) transitions from \((-M,M)\);
\item both rational codes are normalized;
\item the sign bracket holds;
\item the exact width equation \(2^i(b_i-a_i)=2M\) holds.
\end{enumerate}

The base and step proofs use only defining equations and rational order
arithmetic.  The finish proof uses the fuel inequality to derive the requested
width.  These fields instantiate \texttt{TraceObligations};
\texttt{TraceObligations.compile} then produces the embedded implication.  The
archive's Python certificate is an exact regression test for this design, not
a substitute for those object derivations.

\section{Worked example: relational paths without coherent choice}

Use the finite layers
\[
  (0,1,2),\quad(1,2,3),\quad(3,4),\quad(4,5)
\]
and compatibility \(|y-x|\le1\) between adjacent layers.  Backward viability
gives
\[
  (1,2),\quad(2,3),\quad(3,4),\quad(4,5).
\]
The lexicographically first viable path is
\((1,2,3,4)\).  Other paths exist, so a proof architecture demanding a unique
coherent path would create unnecessary work.

Take the observation to be the normalized endpoint \(y/5\).  Across all valid
paths its exact diameter is \(1/5\).  Thus any valid path determines the
observation to tolerance \(1/5\), and a downstream client asking for coarser
precision can ignore the tie entirely.  At finer precision it must request
stronger compatibility or a gap.  This small example is structurally identical
to root matching through collisions, chart gluing, or choosing a carrier map:
backward viability proves extendibility; stable comparison, not coherent
selection, proves the observable.

\section{Lean formalization architecture}
\label{sec:lean}

\subsection{Exact upstream interfaces used}

The supplied modules target the following current definitions:

\begin{itemize}
\item \texttt{Fm.QF} marks quantifier-free formulas;
\item \texttt{Prf Ax Gamma phi} is inductive intuitionistic natural deduction;
\item \texttt{Prf.trans} is simultaneous cut;
\item \texttt{PRA.qfInd} applies QF induction;
\item \texttt{Prf.allE}, \texttt{Prf.exI}, and implication rules perform final
      extraction;
\item \texttt{pra\_iff\_native.mp} transports a completed iPRA derivation to
      the native calculus.
\end{itemize}

The source manifest records the audited commit and SHA of each relevant file.

\subsection{Generic modules}

\texttt{lean/KernelCompiler.lean} contains:

\begin{enumerate}
\item \texttt{checked}: equality of a fixed PR checker to one;
\item \texttt{ProofBuild.nil}, \texttt{step}, \texttt{stepFrom}, and
      \texttt{get};
\item \texttt{invariantAtUnder}, \texttt{loopUnder}, and
      \texttt{loopContract};
\item explicit witness export;
\item stable-observable, boundary, and exact-from-gap rules;
\item native-calculus export.
\end{enumerate}

\texttt{lean/ArithmeticInterface.lean} defines obligation records.  A domain
packet cannot instantiate them without providing actual embedded derivations.
Their \texttt{compile} methods contain only generic logic.

\subsection{Dependency-ordered completion plan}

A complete integration should proceed in the following order.

\begin{enumerate}
\item Compile the generic modules against Lean 4.32.0 and run
      \texttt{\#print axioms}.  Repair API mismatches only.
\item Implement one shared fixed PR coding library for Boolean folds, bounded
      lists/tables, signed integers, rationals, and resource vectors.
\item Prove object-level defining equations and Boolean reflection lemmas.
\item Implement bounded minimum and table folds with QF invariants.
\item Implement a generic fixed trace runner from primitive recursion and fill
      \texttt{TraceObligations} for one benchmark.
\item Formalize backward viability and least viable path over coded tables.
\item Add stable comparison and gap modules.
\item Migrate domain packets one at a time, keeping global bounds as explicit
      boundaries.
\end{enumerate}

This order prevents a large host-level development from drifting away from the
deep embedding.  Every phase ends with a small theorem of type
\texttt{Prf PRAAx}.

\subsection{Expected proof obligations for the cubic benchmark}

The first arithmetic instance should prove:

\begin{enumerate}
\item correctness of rational encoding and comparison;
\item the algebraic identity in \cref{lem:cubic-monotone};
\item initial bracket signs;
\item preservation of sign and width by each branch;
\item correctness of the PR ceiling-logarithm bound;
\item terminal extraction of a width-\(2^{-p}\) bracket.
\end{enumerate}

All are elementary bounded arithmetic.  Completing this benchmark would test
every generic layer without depending on a deep external theorem.

\section{Migrating the existing theorem packets}

For each packet, create a claim table with one row per advertised theorem and
the following columns:

\begin{longtable}{p{0.18\textwidth}p{0.23\textwidth}p{0.23\textwidth}p{0.25\textwidth}}
\toprule
Claim & Finite producer & Embedded verifier & Boundary/gap \\
\midrule
\endhead
Statement consumed by clients & Fixed PR code and explicit resource bound &
Exact \(\Prf\,\PRAAx\) base/step/finish or local DAG proof & Any global theorem,
modulus, separation, or cutoff not generated by the producer \\
\bottomrule
\end{longtable}

Then apply the following triage.

\begin{enumerate}
\item If all four columns are filled, generate the proof schedule and compile.
\item If the producer is missing but a finite witness is supplied externally,
      prove a conditional checker theorem.
\item If only a comparison theorem is available, retain a relational/stable
      observable rather than choosing a representative.
\item If a deep bound is missing, state it as a boundary formula and stop
      claiming unconditional effectivity.
\item If exactness is requested without a gap, weaken the conclusion to the
      strongest certified approximation.
\end{enumerate}

This procedure directly addresses the difficulty reports.  It does not discard
ambitious packets: it separates completed finite mathematics from the exact
open lemma blocking an unconditional object proof.

\section{Limitations and nonclaims}

The following limitations are mathematical, not merely engineering details.

\begin{itemize}
\item The framework does not prove effective forms of the major global
      arithmetic theorems whose known arguments do not provide the required
      bounds.
\item Finite viability eliminates bounded choice, not arbitrary choice.
\item Stable comparison gives an observable only at the proved tolerance.
      Exact labels still require a gap.
\item Fixed-program specialization avoids, rather than solves, the universal
      evaluator obstruction.
\item A human proof that an algorithm is primitive recursive is not yet its
      object-iPRA correctness proof.  Defining equations, QF invariants, and
      extraction derivations must be supplied.
\item The included Lean source has not been compiled in the artifact runtime.
      Its lack of placeholders is weaker than successful typechecking and
      axiom auditing.
\end{itemize}

These limits make the package more useful: every future claim has a clear place
in the architecture and a precise completion criterion.

\section{Conclusion}

The ProveIt iPRA corpus has converged on a strong insight: advanced mathematics
must be reorganized around finite observations, explicit resources, bounded
certificates, and quantitative stability.  The remaining systemic problem is
not another semantic abstraction.  It is the final kernel-facing compiler that
turns local proof objects into an actual theorem of the deep embedding.

The proof-carrying schedule developed here supplies that compiler for every
fixed finite acyclic contract graph.  Quantifier-free loop reflection supplies
its iterative nodes.  Backward viability removes hidden finite choice.  Stable
observables remove unnecessary coherent representatives.  Gap rules prevent
invalid exactness claims.  Boundary formulas expose deep global mathematics
instead of burying it inside an alleged algorithm.

The immediate research program is therefore concrete: typecheck the generic
modules, complete one shared coding library, formalize the cubic benchmark, and
migrate existing packets by filling explicit obligation records.  Each step
produces reusable \(\Prf\,\PRAAx\) theorems and converts an informal difficulty
report into a sharply located mathematical task.

\appendix

\section{Theorem dependency graph}

\begin{center}
\begin{tabular}{p{0.23\textwidth}p{0.67\textwidth}}
\toprule
Result & Dependencies \\
\midrule
Proof-carrying step & \texttt{Prf.trans}, hypothesis rule, list membership \\
Finite DAG compiler & Proof-carrying step, finite metalevel induction \\
Loop compiler & \texttt{PRA.qfInd}, universal elimination, cut, implication introduction \\
Bounded least search & Fixed PR coding, QF induction, Boolean folds \\
Viability correctness & Bounded least search, backward finite induction \\
Stable observable & Implication introduction and elimination \\
Gap decision & Rational triangle inequalities and explicit separation \\
Boundary factorization & Implication introduction and elimination \\
Cubic certificate & Exact rational arithmetic, monotonicity identity, loop compiler \\
\bottomrule
\end{tabular}
\end{center}

\section{Object-formula schemas}

A practical library should standardize the following formulas, all with fixed
PR symbols and explicit arities.

\begin{align*}
  \mathsf{Checked}_b(\vec x)
    &:\equiv b(\vec x)=1,\\
  \mathsf{Graph}_f(\vec x,y)
    &:\equiv f(\vec x)=y,\\
  \mathsf{TraceInv}(i,s,\vec x)
    &:\equiv \mathsf{run}(i,\vec x)=s\ \wedge\
      \mathsf{allLocalChecks}(i,s,\vec x)=1,\\
  \mathsf{Viable}(i,x,T)
    &:\equiv \mathsf{tableGet}(T,i,x)=1,\\
  \mathsf{Gap}(q,\delta)
    &:\equiv q=0\ \vee\ \delta\le |q|,\\
  \mathsf{Stable}_p(c_1,c_2)
    &:\equiv \dist(\obs(c_1),\obs(c_2))\le 2^{-p}.
\end{align*}

Bounded conjunctions and disjunctions are represented by checker terms rather
than object quantifiers, preserving the QF status of induction invariants.

\section{Source and verification status}

The source baseline is ProveIt main commit
\texttt{9df5b64f82b891640b2154bd14dcee251a25d37f}.  The exact file hashes used for
the formal interface are recorded in
\texttt{metadata/source\_manifest.json}.  The archive contains:

\begin{itemize}
\item this LaTeX source and compiled PDF;
\item placeholder-free generic Lean source, not typechecked in the build
      runtime;
\item a generated Lean schedule example;
\item exact Python reference algorithms and unit tests;
\item JSON certificates and a reproducible verification report;
\item a theorem ledger, difficulty matrix, formalization roadmap, and SHA-256
      manifest.
\end{itemize}

\begin{thebibliography}{9}
\bibitem{ProveItPRA}
V. Reshetnikov,
\emph{ProveIt: Logic/PrimitiveRecursiveArithmetic}, current repository tree,
\url{https://github.com/VladimirReshetnikov/ProveIt/tree/main/Logic/PrimitiveRecursiveArithmetic}.

\bibitem{ConciseIPRA}
V. Reshetnikov et al.,
\emph{ConciseIPRA}, deep-embedded Lean formalization in the ProveIt repository,
audited at the commit recorded in the source manifest.

\bibitem{Kleene}
S. C. Kleene,
\emph{Introduction to Metamathematics},
North-Holland, 1952.

\bibitem{TroelstraVanDalen}
A. S. Troelstra and D. van Dalen,
\emph{Constructivism in Mathematics}, Volumes I--II,
North-Holland, 1988.
\end{thebibliography}

\end{document}
